%%%%%%%%%%%%%%%%%%%%%%%%%%%%%%%%%%%%%%%%%%%%%%%%%%%%%%%%%%%%
%% CHAPTER 2 -- How Machines Learn from Their Mistakes
%%%%%%%%%%%%%%%%%%%%%%%%%%%%%%%%%%%%%%%%%%%%%%%%%%%%%%%%%%%%

\chapter{How Machines Learn from Their Mistakes}
\label{chap:optimization}
\label{ntg:ch2}

\scenelabel
\begin{scenestrip}
\sceneitem{The problem}
\sceneitem{How we got here}
\sceneitem{Where we stand}
\sceneitem{What goes wrong}
\end{scenestrip}

\paragraph{The problem.} Chapter~\ref{chap:stat-foundations} told us \emph{what} we want a language model to do:
look at the words so far and produce a sensible weighting over
the words that might come next. It said almost nothing about
\emph{how} the model came to be able to do that. We now face
that question. The model has, on the small side, a hundred
million internal numbers --- the parameters --- and on the
large side a trillion. Each of those numbers needs to be set
to a value that makes the model's predictions, on average,
correct. There is no clever shortcut. There is no formula one
can solve in closed form. There is, instead, a procedure: try
something, measure how wrong it was, change the numbers a
little, try again, and repeat for months on end.

The idea seems almost too simple to be the basis of an entire
field. It is. The remarkable thing is that the procedure works
at all. The slightly less remarkable thing is what it costs in
electricity, in chips, in human attention, and in the design
choices that surround it.

\paragraph{How we got here.} The phrase "change the numbers a little"
conceals a question. A large language model has hundreds of millions
to hundreds of billions of parameters. Searching through all possible
combinations at random would take longer than the age of the universe.
You need to know \emph{which} numbers to change, and \emph{in which
direction}.

Here is the key insight. For any particular setting of all the
parameters, you can compute the wrongness score. That makes the
wrongness a function of the parameter settings: a single number
that goes up or down as you move through the space of all possible
combinations. Imagine --- as a mental shorthand, since the real
space has as many dimensions as the model has parameters --- that
there are only two parameters, and you draw them as two axes on a
map. At every point on the map, the height is the wrongness for
that combination. Training is the problem of finding the lowest
valley on this map.

Now suppose you are standing somewhere on that landscape and can
feel only the slope of the ground immediately beneath your feet.
The slope is informative: if the ground runs downhill to your left,
take a step left. The same logic holds in any number of dimensions.
For each parameter, you can ask whether nudging it very slightly
upward or downward would make the wrongness increase or decrease,
and by how much. This rate of change --- the slope of the wrongness
with respect to each parameter, computed at the current setting ---
is a list of numbers, one per parameter. Mathematicians call this
list the \emph{gradient}. It tells you, at your current position in
the landscape, which direction runs downhill. Take a small step in
that direction, recompute the gradient, and repeat. This is the
core of the whole procedure.

The mathematics of descending a function by following its slope
is older than electricity. In 1847 the French mathematician
Augustin-Louis Cauchy described, in a short note to the French
Academy of Sciences, the method we now call \emph{gradient
descent}. The idea is one a hiker would recognise. If you want
to reach the bottom of a valley in a fog so thick that you can
see only your own feet, the best you can do is feel the slope
of the ground beneath you, take a small step downhill, and
repeat. You may not reach the lowest possible point in the
landscape. You will, if the landscape is reasonable, reach a
point lower than where you started. Cauchy's note was about
solving small algebraic problems. The idea would wait a
century before becoming the dominant tool of artificial
intelligence.

The trick that bridged Cauchy's idea to modern machine learning
came in 1951, in a paper by the American statisticians Herbert
Robbins and Sutton Monro. They observed that, in many practical
problems, computing the exact slope of the function you were
trying to descend was impossible --- you could only get a
\emph{noisy estimate} from a small sample of data. They proved
that, under reasonable conditions, descending in the direction
of the noisy estimate would still get you somewhere useful, as
long as you took progressively smaller steps. They called this
\emph{stochastic approximation}. We now call it stochastic
gradient descent, and it is the engine of every neural network
training run that has ever been performed.

The other trick we needed was a way to compute the slope of a
deep network at all. The chain rule of calculus answers this in
principle. The implementation is called \emph{back-propagation},
and the name is fair: it computes the slope by walking backwards
through the network from output to input. Paul Werbos, in his
1974 Harvard PhD thesis, wrote down a recognisable version of
the algorithm for neural networks. The field did not notice. It
was rediscovered, polished, and brought to general attention by
David Rumelhart, Geoffrey Hinton, and Ronald Williams in their
1986 paper in \emph{Nature}. Without back-propagation, a network
of more than a few layers could not be trained at all. With it,
a network of a hundred layers could be trained almost as easily
as a network of one. The remainder of the field's history is
what people did with the algorithm Werbos had written down and
the world had forgotten for twelve years.

For thirty years after that the field argued about which
flavour of gradient descent was best. Momentum (which lets the
descent build up speed in consistent directions, like a heavy
ball rolling down a slope). RMSProp, AdaGrad, AdaDelta. In 2014
two graduate students, Diederik Kingma and Jimmy Ba, combined
several of the existing ideas into a recipe they called
\emph{Adam}. With minor refinements --- AdamW from Loshchilov
and Hutter in 2017 fixed a subtle issue with how the algorithm
handled a regularisation trick called weight decay --- this is
still the optimiser used to train almost every Large Language
Model in production. Choosing an optimiser is no longer a
research question. The research questions are now elsewhere.

The most consequential of those questions is: \emph{how much
of what should we spend?} Jared Kaplan and his coauthors at
OpenAI, in 2020, observed that the loss --- the surprise number
from Chapter~\ref{chap:stat-foundations} --- decreases as a slow, predictable function of
three quantities: the number of parameters in the model, the
number of words in the training corpus, and the total amount of
arithmetic spent. The relationship is what mathematicians call
a \emph{power law}, which means that doubling any of the three
quantities buys you a fixed proportional improvement in the
loss. Power laws are not exciting. They predict that progress
will be steady and expensive rather than sudden and free.

Two years later Jordan Hoffmann and her colleagues at DeepMind
ran a more careful version of the same study and reached a
sharper conclusion. They argued that most existing large
models had been trained with too few words for the size they
were --- that, given a fixed compute budget, you got a better
model by training a smaller network on more text than by
training a larger network on less. They named the recipe after
their flagship model: \emph{Chinchilla}. The Chinchilla
argument changed how compute is allocated across the industry,
and it is the reason the headline parameter count of new
models has, in the past three years, mostly stopped going up.

\paragraph{Where we stand.} A modern training recipe is a long sequence of small choices,
each of which has been worked out by trial and error over a
decade. Choose AdamW. Pick a peak learning rate (a number that
controls how big each step downhill is) and start it at zero,
ramping it up slowly over the first few thousand steps and then
gradually shrinking it back down over the rest of the run. Use
mixed precision arithmetic --- half precision for most of the
work, full precision where it matters for stability. Clip
unusually large gradients so a single bad batch cannot wreck
months of training. Pick a model size and a corpus size in the
proportion the Chinchilla paper recommended. Run on a cluster of
thousands of high-end accelerator chips. Watch the loss curve
the way an air traffic controller watches a screen. Recover
from the failures, of which there will be many. After several
weeks to several months you will have a trained model.

The ``thousands of high-end accelerator chips'' is not a
rhetorical flourish. The training of a frontier model in
2024 used somewhere between roughly tens of thousands of
H100-class GPUs running for several weeks. The marginal cost
of electricity, hardware depreciation, and operational support
ran into the tens to hundreds of millions of dollars. The
total upfront capital expenditure on the data centres that
host these training runs ran into the tens of billions. These
numbers are accurate as of this writing; they will be larger
when you read this; they are themselves the subject of
considerable industry secrecy and may be off by a factor of
two in either direction. The qualitative point is that
training a frontier model is a financial and industrial
undertaking on the scale of building a new aircraft.

\paragraph{What goes wrong.} The dramatic failures of training are visible. A loss curve that
suddenly spikes upward and never recovers. A gradient that
explodes after fifty thousand steps and turns the rest of the
run into nonsense. A divergence at the end of week three that
turns four million dollars of cluster time into nothing. These
failures look bad and they cost money, but they are the easy
ones. The engineer who notices them can usually diagnose them:
a learning rate set too high, a batch of training data with an
outlier in it, a precision setting too aggressive. They show up
on the dashboard and they get fixed.

The harder failures are the quiet ones. A model that trains to
a respectable score on the held-out evaluation set and yet
behaves badly on the slightly different distribution of inputs
it will actually face when users send it questions. A model
that has, by accident, memorised passages of its training
corpus rather than learned the patterns the corpus exhibits ---
which is fine until a user asks the right question and the
model returns a verbatim copy of someone's copyrighted novel.
A model whose evaluation scores improved relative to the
previous generation because the evaluation set was, unbeknownst
to the engineers, contained in the training data. A model that
appears to have a useful behaviour in testing but which only
exhibits the behaviour because the test prompts were drawn from
the same distribution as the data the model was trained on,
and that fails when the prompts are even slightly different.

These failures are not failures of the optimiser. The optimiser
did exactly what it was asked to do: it reduced the loss. The
failures are mismatches between what the loss measures and what
the deployed system actually needs to do. We will see this
pattern repeat throughout the rest of the book under different
names: in Chapter~\ref{chap:pretraining} it will be the gap between pretraining and
deployment, in Chapter~\ref{chap:alignment} it will be the alignment tax, in
Chapter~\ref{chap:systems} it will be evaluation contamination, in Chapter~\ref{chap:error-model} it
will be the structure-meaning-context-groundedness-environment
ladder.

\section{The idea, in plain language}

\subsection*{The hiker in the fog}

The single most useful image for understanding modern machine
learning is the hiker in the fog. The landscape is the loss
surface --- a multidimensional terrain in which every point
corresponds to one possible setting of the model's parameters,
and the height of the terrain at that point is how badly the
model performs at the values. The hiker's job is to find a low
spot. The fog is so thick that the hiker can only feel the
slope under their feet. They take a step in the steepest
downhill direction. They feel the new slope. They take another
step. After enough steps, they reach a low point. They have no
guarantee it is the lowest. They have only the guarantee that
they are lower than where they started.

Several details of the picture matter for what follows.

First, the dimensionality. The hiker is not walking in three
dimensions. For a Large Language Model, the hiker is walking
in a space with billions of dimensions, one per parameter.
Visualisation fails almost immediately, and so does intuition
based on hill climbing in everyday landscapes. What we know
about the actual landscape is mostly what we have observed
empirically: it has many local low points, most of them are
roughly as good as one another, and the path between them is
mostly downhill. Why this should be true is not, at the time
of writing, well understood theoretically. It is one of the
several places where the field is operating on empirical
observation rather than principle.

Second, the noise. The hiker does not get to feel the true
slope. The slope they feel is computed from a small sample of
training examples, and it is therefore noisy --- a different
sample would produce a slightly different feel. The training
of a Large Language Model takes the noisiness as a feature
rather than a bug. The noise prevents the hiker from getting
stuck on small bumps in the landscape, the way a precise step
might. The right metaphor is a slightly drunk hiker, whose
small random missteps occasionally jiggle them out of holes
they would otherwise sit in forever.

Third, the step size. Each step is small --- much smaller than
the diameter of the bumps and basins on the surface. If the
steps are too large, the hiker overshoots and oscillates. If
they are too small, the hiker takes forever to get anywhere.
The ``learning rate'' is the engineering name for the step
size, and most of the design of a modern training schedule is
about how to vary it across the run --- larger early, when the
hiker is far from any minimum and any direction is likely to
help; smaller late, when the hiker is close enough to a
minimum that overshoot would be catastrophic.

\subsection*{Why scaling sometimes works}

The Kaplan and Chinchilla papers are worth understanding in
their qualitative form even if you never look at their
equations.

The Kaplan paper observed that the loss --- the surprise number
the training is trying to reduce --- decreases predictably as
you spend more on it. Specifically, the loss falls along a
curve whose shape is the same regardless of which of the three
spending quantities (parameter count, dataset size, total
compute) you change. The curve is one of diminishing returns:
each doubling of compute buys a smaller proportional
improvement in loss than the doubling before. But the
improvements do not stop. They continue, predictably, for as
long as anyone has been able to measure them.

This is not the same as saying that every capability of the
model improves smoothly. Some capabilities --- the ones that
require chaining several steps of reasoning together, for
example --- appear suddenly at certain scales, in a way that
looks discontinuous. The technical name for this is
\emph{emergence}, and there is a continuing argument in the
literature about whether it is a real phenomenon or an artefact
of how the capabilities are being measured. We will return to
this in Chapter~\ref{chap:pretraining}.

The Chinchilla paper sharpened the Kaplan picture by pointing
out that, given a fixed budget, you do better by training a
smaller model on more data than by training a bigger model on
less. The argument has held up well empirically. It is the
reason a typical 2024 model is, by parameter count, smaller
than the 2020 GPT-3 (which had 175 billion parameters), and yet
better than it on every benchmark anyone has run. The Chinchilla
ratio is the rough rule of thumb that, for a given size of
model, you want to train on roughly twenty times as many words
as the model has parameters.

\subsection*{Why scaling sometimes does not work}

The Kaplan and Chinchilla curves describe the gross behaviour
of training loss as a function of resources. They do not
describe everything that matters. In particular, they do not
describe deployment behaviour, which is what users care about.

A model that has a lower training loss than another model is
better, on average, at predicting the next word in its training
distribution. This usually translates into better deployment
behaviour, but not always. The translation is mediated by:

\begin{itemize}
  \item How well the training distribution matches the
        deployment distribution. If the training corpus is
        mostly informal web text and the deployment is asking
        the model to draft contracts, the loss improvement may
        not transfer.
  \item How well the alignment work of Chapter~\ref{chap:alignment} has been done.
        A model that is perfectly trained but unhelpfully
        evasive at deployment is not actually useful.
  \item How well the surrounding system --- the retrieval, the
        tool calls, the prompt template --- exploits the model's
        capability. A small well-equipped model can outperform
        a large under-equipped one.
\end{itemize}

The honest summary is that scaling helps but is not a strategy
on its own. The gap between a research model that scores well
on internal benchmarks and a product that delights users is not,
in our experience, predictable from the benchmarks. We will
return to this in Chapter~\ref{chap:systems} when we discuss evaluation
harnesses, and again in Chapter~\ref{chap:governance} when we discuss what
``deployment-ready'' actually means.

\section{The engineering consequence}

The fact that training is a procedure rather than an algorithm
has several consequences worth holding in mind for any reader
trying to reason about a real LLM product.

\paragraph{Training is a slow, expensive, opinionated process.}
Every training run involves hundreds of decisions about
hyperparameters, data composition, and evaluation. Most of those
decisions are not documented in any meaningful way. When a
vendor tells you that ``model X is better than model Y'', they
are reporting the outcome of a specific recipe. A different
recipe with the same starting materials would produce a
different model, and the vendor will, in most cases, not tell
you which choices in the recipe mattered most.

\paragraph{Compute, data, and parameters are jointly the budget.}
A small organisation cannot, by training a large model, beat a
large organisation that has trained a smaller one on much more
data. The Chinchilla insight is that, for any compute budget,
there is an efficient frontier of (model size, training data
size) pairs, and most attempts to beat the frontier by spending
more on one of the three quantities than on the others are
wasteful. Reasoning about whose model is better, in this
context, is mostly reasoning about who got closer to the
frontier with the budget they had.

\paragraph{The training cost is not the deployment cost.} A
model that cost a hundred million dollars to train may cost
thousandths of a cent to use for a single query. The mismatch is
deliberate: training is amortised across hundreds of millions
of queries. It also has a perverse consequence: an organisation
that trained a frontier model has no economic reason to retrain
it from scratch when a better recipe appears. They have an
economic reason to fine-tune (Chapter~\ref{chap:adaptation}), to add retrieval
(Chapter~\ref{chap:systems}), to align (Chapter~\ref{chap:alignment}). Training a brand new
foundation model is a bet that is made only by the small set of
organisations that can afford to make it, and it is made
roughly once a year per such organisation.

\paragraph{Training has its own failure modes that matter
operationally.} A model that has been trained on data of poor
quality cannot be repaired by post-hoc alignment. A model whose
training data was contaminated with the test set will look
better than it is until it meets reality. A model trained
without sufficient diversity in its corpus will be quietly
worse on populations underrepresented in that corpus. None of
these failures are visible at the API layer. All of them are
deeply consequential for whoever has to defend the deployed
behaviour to a regulator, an auditor, or a sceptical user.

\section{What goes wrong, and why}

Several characteristic failure modes trace back to the training
process described in this chapter.

\paragraph{Memorisation.} Because the training process is
shown each example in its corpus repeatedly, and because a
model with a billion parameters has more than enough capacity
to store many of those examples verbatim, modern models reliably
memorise some fraction of their training data. The fraction is
small in proportional terms but large in absolute terms.
Memorised content is recoverable from the model under the right
prompt, which has implications for copyright (Chapter~\ref{chap:governance}), for
privacy (a model trained on a corpus that contained leaked
personal data may surface that data on request), and for
benchmark integrity (a model that memorised the test set will
look better than its actual capability). The mechanism is the
training process doing what it was asked to do; the consequence
is one the field is still working out how to manage.

\paragraph{Distribution shift.} The model was trained on data
drawn from a particular distribution: the public web up to a
certain date, certain books, certain code repositories. The
model is then asked to perform on inputs drawn from a different
distribution: queries from real users in the present day, often
about subjects that did not exist when the training data was
collected, or framed in ways that do not appear in the training
data. The gap between training distribution and deployment
distribution is the single largest source of unpredictable
behaviour in the field. We will see in Chapter~\ref{chap:systems} how retrieval
helps to close part of the gap, in Chapter~\ref{chap:alignment} how alignment
helps to close another part, and in Chapter~\ref{chap:error-model} how, even with
both interventions, a residue remains.

\paragraph{Compute as a binding constraint.} A widespread
public misunderstanding holds that the field is bottlenecked by
algorithmic ideas. This is, at the time of writing, not the
view of most researchers in the field. The field is bottlenecked
by compute, by the supply of high-quality training data, and by
the supply of trained engineers. Algorithmic improvements
matter, but their effect on the deployed product is, in
practice, a multiplier of an order of magnitude or so on top of
the gross effect of more compute and better data. This has
implications for how to read industry announcements: most
``breakthroughs'' are points along the scaling curve rather
than departures from it.

\paragraph{Training data choices are quietly load-bearing.}
Two models with identical architectures and identical
parameter counts can behave very differently if their training
corpora differ. Curating the corpus --- which texts to include,
which to exclude, in what proportions, with what filtering
applied --- is one of the most important and least visible
parts of building a modern model. The decisions are usually
not documented publicly, often for competitive reasons and
sometimes for reasons of legal exposure. The result is that
even a careful evaluator can have difficulty saying \emph{why}
one model is better than another, because the answer is often
``the training data was different in ways the vendor will not
disclose''.

\section*{Bridges}
\addcontentsline{toc}{section}{Bridges}

This chapter depended only on Chapter~\ref{chap:stat-foundations}: we are reducing the
``surprise'' number Chapter~\ref{chap:stat-foundations} introduced. The hiker-in-the-fog
picture and the Chinchilla scaling argument reappear in
Chapter~\ref{chap:pretraining} (where we apply them to actual pretraining), in
Chapter~\ref{chap:efficiency} (where we discuss what to do when the budget is
binding), and in Chapter~\ref{chap:adaptation} (where we discuss adaptation as the
art of getting useful model behaviour without paying the full
training cost). The notion that the loss is a proxy for the
deployment property the user actually cares about reappears
throughout the rest of the book, and is the central organising
idea of Chapter~\ref{chap:error-model}. The next chapter --- \emph{Networks of
Simple Decisions} --- explains, in similarly plain language,
what is actually inside the network that the procedure of this
chapter is training.
